\section{The data hiding game}\label{sec:data-hiding-layer}

In this section, we introduce a new, simple game called the \emph{data hiding game}.
This game assumes two $(k,n,q)$-semiregister provers with a shared state $\ket{r_1} \cdots \ket{r_k} \ket{\mathrm{aux}}$.
The goal is to test that a given measurement $\{M_a^x\}_a$ acts as the identity on the $k$-th register.

\begin{definition}\label{def:data-hide-def}
Let $x = (x_1, x_2)$ with $x_1 = (W_1, \ldots, W_k)$, and suppose $W_k = \hideq$.
Then the \emph{data hiding game}~$\game_{\mathrm{hide}} := \game_{\mathrm{hide}}(x)$ is given by~\Cref{fig:hide-one}.
It has the following parameters:
\begin{equation*}
\qtime{\game_{\mathrm{hide}}},
\qlength{\game_{\mathrm{hide}}} = O(|x|),
\quad
\atime{\game_{\mathrm{hide}}},
\alength{\game_{\mathrm{hide}}} = O(\textstyle{\sum_i} n_i \log(q_i) + \ell).
\end{equation*}
Here $\ell$ is the maximum of $|a_2|, |a_2'|$ over all answers $a_2$ and $a_2'$ given by the provers.
\end{definition}

{
\floatstyle{boxed} 
\restylefloat{figure}
\begin{figure}[htbp]
Draw $\bW \sim \{X, Z\}$. Set $\bx' = (\bx_1', x_2)$, where $\bx_1' = (W_1, \ldots, W_{k-1}, \bW)$.
Flip an unbiased coin $\bb \sim \{0, 1\}$. Distribute the questions as follows:
\begin{itemize}
\item[$\circ$] Player~$\bb$: give $x$; receive $(\ba_1, \ba_2)$.
\item[$\circ$] Player~$\overline{\bb}$: give $\bx'$; receive $(\ba_1', \ba_2')$.
\end{itemize}
Accept if $\ba_2 = \ba_2'$.
\caption{The game $\game_{\mathrm{hide}}(x)$, with input $x = (x_1, x_2)$ \label{fig:hide-one}}
\end{figure}
}

For a measurement $\{M_a\}_a$ which operates on multiple subsystems,
it will be convenient to define a version of the measurement in which one of the subsystems is ``hidden".

\begin{notation}
Let $M$ be a matrix which operates on $\mathcal{H}_1 \otimes \cdots \otimes \mathcal{H}_k \otimes \mathcal{H}_{\mathrm{aux}}$, and let $i \in [k]$.
Define the notation
\begin{equation*}
\mathrm{hide}_i(M)
:= \frac{1}{\tr(I_i)} \cdot I_i \otimes \tr_i(M).
\end{equation*}
\label{not:hide}
\end{notation}
If $\{M_a\}_a$ is a measurement, then so is $\{\hide{k}{M_a}\}_a$ (though it may not be projective, even if $\{M_a\}_a$ is).
Our main result regarding the data hiding game is that passing it with high probability certifies that $\{M_a\}_a$ is close to $\{\hide{k}{M_a}\}$.

\begin{theorem}\label{thm:data-hiding-game}
Let $x$ be as in \Cref{def:data-hide-def}.
\begin{itemize}
\item[$\circ$] \textbf{Completeness:}
Let $\calS_{\mathrm{partial}} = (\psi, M^x)$ be a partial $(k,n,q)$-register strategy
			which is also a real commuting EPR strategy.
			Then there is a $(k,n,q)$-register
                        strategy~$\calS$ extending
                        $\calS_{\mathrm{partial}}$ which is also a
                        real commuting EPR strategy
			such that $\valstrat{\game_{\mathrm{hide}}}{\calS} = 1$.
\item[$\circ$] \textbf{Soundness:} 
Let $\calS = (\psi, M)$ be a projective $(k, n, q)$-semiregister strategy such that $\valstrat{\game_{\mathrm{hide}}}{\calS} \geq 1-\epsilon$. Then
\begin{equation*}
(M_a^x)_{\reg{Alice}} \otimes I_{\reg{Bob}} \approx_\eps (\hide{k}{M_a^x})_{\reg{Alice}} \otimes I_{\reg{Bob}}
\end{equation*}
on the singleton distribution on input~$x$.
\end{itemize}
\end{theorem}

This section is organized as follows:
in \Cref{sec:pauli-twirl} we introduce the \emph{Pauli twirl},
and in \Cref{sec:data-hiding-game} we use it to prove \Cref{thm:data-hiding-game}.
Finally, in \Cref{sec:compile-sec}, we design our compiler from layer-two to layer-one.
This last step is essentially standard and is included for completeness.

\subsection{Some facts about the Pauli twirl}\label{sec:pauli-twirl}

\newcommand{\push}{\nu}
\newcommand{\bpush}{\bnu}
\newcommand{\twirlset}{S_{\mathrm{twirl}}}
\newcommand{\twirldist}{\calD_{\mathrm{twirl}}}

\begin{definition}
  The \emph{Pauli twirl} $\twirl: \calB((\mathbb{C}^q)^{\ot n})
  \to \calB((\mathbb{C}^q)^{\ot n})$ is the linear operator
  \[ \twirl(A) := \E_{\bu, \bu' \sim \F_q^n}\left[
      X(\bu) Z(\bu') \cdot A \cdot Z(-\bu') X(-\bu) \right]. \]
\end{definition}

\begin{proposition}\label{claim:twirled_pauli}
  Let $P$ be a Pauli matrix on $n$ qudits of dimension $q$. Then
$
\twirl(P) =  P
$
if $P$ is a multiple of the identity,
and otherwise $\twirl(P) = 0$.
\end{proposition}

\begin{proof}
The case when~$P$ is a multiple of the identity follows from the definition.
Otherwise, we can write $P = \omega^z X(a) Z(b)$, where at least one of~$a$ and~$b$ is nonzero. Then
\begin{align*}
     \twirl(P)
&=  \E_{\bu, \bu'} \left[ X(\bu) Z(\bu') \cdot P \cdot
		Z(-\bu') X(-\bu) \right]  \\
&= \omega^z \E_{\bu, \bu'} \left[ X(\bu) Z(\bu') \cdot X(a) Z(b) \cdot
		Z(-\bu') X(-\bu) \right].
\end{align*}
By the Pauli~$X$ and~$Z$ commutation relations (\Cref{eq:commutation-relations}),
this rearranges to
\begin{equation*}
 \omega^z \E_{\bu, \bu'} \left[\omega^{\tr[\bu' \cdot a - \bu \cdot b]}\right] \cdot X(a) Z(b)
=  \E_{\bu, \bu'} \left[ \omega^{\tr[\bu' \cdot a - \bu \cdot b]} \right] \cdot P
    =  \E_{\bu'} \big[ \omega^{\tr[\bu' \cdot a]}\big] \cdot \E_{\bu}\big[\omega^{- \tr[\bu \cdot b]} \big] \cdot P
    = 0.
\end{equation*}
Here the last step uses \Cref{fact:averages-to-zero} and the fact that at least one of~$a$ or~$b$ is nonzero.
\end{proof}

In the next couple of sections,
we will consider the effects of applying the Pauli twirl to our measurements.
For convenience, we will ``group" our state into two parts: $\ket{\psi_1} = \ket{r_k}$ 
is the subsystem we want to hide,
and $\ket{\psi_2} = \ket{r_1} \cdots \ket{r_{k-1}} \ket{\mathrm{aux}}$ is the remaining part of this state.
In this way, we can consider our measurements as operating on the bipartite state $\ket{\psi_1} \ket{\psi_2}$.

\begin{proposition}\label{prop:hidden-twirl}
Let $\{M_a\}$ be a measurement on the state $\ket{\psi} = \ket{\psi_1} \ket{\psi_2}$.
Then
\begin{equation*}
(\twirl_1\otimes \id_{2})[M_a]
=
\hide{1}{M_a},
\end{equation*}
where $\id_{2}$ is the identity superoperator applied to the second register.
\end{proposition}
\begin{proof}
Let $P_J$ be the elements of the Pauli group on $n$ qudits of dimension~$q$, with $P_0 = I$.
Because these form a basis for the set of matrices, we can write
\begin{equation*}
M_a = \sum_{J} P_J \otimes M_{a, J},
\end{equation*}
where the $M_{a, J}$'s are matrices acting on the auxiliary register. 
Using \Cref{claim:twirled_pauli},
\begin{equation*}
(\twirl_1\otimes \id_{2})[M_a]
= \sum_{J}  \twirl(P_J) \otimes M_{a,J}
= P_0 \otimes M_{a,0}
= I \otimes M_{a,0}.
\end{equation*}
On the other hand, because $P_J$ is traceless unless $J=0$ (i.e.\ $P_J$ is the identity),
\begin{equation*}
\hide{1}{M_a}
= \sum_{J} \hide{1}{P_J \otimes M_{a,J}} 
= \sum_{J} \frac{1}{q^n} \cdot I \otimes \tr_1(P_J \otimes M_{a,J})
= I \otimes M_{a, 0}.
\end{equation*}
These two are equal, completing the proof.
\end{proof}

\subsection{Hiding a single coordinate}\label{sec:data-hiding-game}

In this section, we prove \Cref{thm:data-hiding-game}. Prior to doing so, we prove a couple of technical lemmas.
The first shows that a measurement which approximately commutes with the Pauli measurements
also approximately commutes with the Pauli observables.

\begin{lemma}\label{lem:proj-to-obs}
Let $W \in \{X, Z\}$.
Suppose $\{M_{a}\}$ is a measurement on the state $\ket{\psi} = \ket{\mathrm{EPR}_q^n} \ket{\psi_2}$ for which
\begin{equation*}
(M_{a} \cdot (\tau^{W}_{u} \otimes I_{\reg{2}}))_{\reg{Alice}} \otimes I_{\reg{Bob}}
\approx_{\delta} ((\tau^{W}_{u} \otimes I_{\reg{2}}) \cdot M_{a})_{\reg{Alice}} \otimes I_{\reg{Bob}}.
\end{equation*}
Then the statement
\begin{equation*}
(M_{a} \cdot (W(u) \otimes I_{\reg{2}}))_{\reg{Alice}} \otimes I_{\reg{Bob}}
\approx_{\delta} ((W(u) \otimes I_{\reg{2}}) \cdot M_{a})_{\reg{Alice}} \otimes I_{\reg{Bob}}
\end{equation*}
holds with respect to the uniform distribution on $\bu \in \F_q^n$.
\end{lemma}
\begin{proof}
Our goal is to bound
\begin{equation}\label{eq:gonna-bound-this-foist}
\E_{\bu} \sum_{a} \Vert (M_{a} \cdot (W(\bu) \otimes I_{\reg{2}}) - (W(\bu) \otimes I_{\reg{2}}) \cdot M_{a}) \otimes I \ket{\psi} \Vert^2.
\end{equation}
by $\delta$.
To do so, for a fixed~$u$ we introduce the notation
\begin{align}
\Delta_{a}^{u}
&:= M_{a} \cdot (W(u) \otimes I_{\reg{2}}) - (W(u) \otimes I_{\reg{2}}) \cdot M_{a}\nonumber\\
&= \sum_{v \in \F_q^n} \omega^{\tr[u \cdot v]} (\underbrace{M_{a}
	\cdot (\tau^{W}_{v} \otimes I_{\reg{2}})- (\tau^{W}_v \otimes I_{\reg{2}}) \cdot M_{a}}_{\Delta_{a,v}}).\label{eq:pauli-obs-commutator-foist}
\end{align}
We record the following identity, which follows from~\Cref{eq:pauli-obs-commutator-foist}:
\begin{equation*}
\E_{\bu} (\Delta_{a}^{\bu})^\dagger\cdot \Delta_{a}^{\bu}
= \E_{\bu} \sum_{v, v' \in \F_q^n} \omega^{\tr[\bu \cdot (v' - v)]} (\Delta_{a,v})^\dagger \Delta_{a,v'}
= \sum_{v \in \F_q^n} (\Delta_{a,v})^\dagger \Delta_{a,v}.
\end{equation*}
As a result,
\begin{align*}
\eqref{eq:gonna-bound-this-foist}=
\E_{\bu} \sum_{a}\Vert (\Delta_{a}^{\bu} \otimes I) \ket{\psi}\Vert^2
&= \E_{\bu} \sum_{a} \bra{\psi} (\Delta_{a}^{\bu})^\dagger \Delta_{a}^{\bu} \otimes I \ket{\psi}\\
&=  \sum_{a} \sum_{v \in \F_q^n} \bra{\psi}(\Delta_{a,v})^\dagger \Delta_{a,v} \otimes I \ket{\psi}\\
&= \sum_{a, v} \Vert \Delta_{a,v} \otimes I \ket{\psi} \Vert^2.
\end{align*}
But this is at most $O(\delta)$, by assumption. This completes the proof.
\end{proof}

The next technical lemma shows that a measurement which approximately commutes with products of $X$ and $Z$ observables
is approximately equal to its own Pauli twirl.

\begin{lemma}\label{lem:commute-to-twirl}
Consider the distribution~$\calD$ on pairs $(\bu, \bu')$, where $\bu, \bu' \sim \F_q^n$.
Suppose $\{M_{a}\}$ is a measurement on the state $\ket{\psi} = \ket{\mathrm{EPR}_q^n} \ket{\psi_2}$ for which
\begin{equation*}
((Z(u') X(u)\otimes I_{\reg{2}}) \cdot M_{a}) \otimes I_{\reg{Bob}}
\approx_{\delta} (M_{a} \cdot (Z(u') X(u)\otimes I_{\reg{2}})) \otimes I_{\reg{Bob}}.
\end{equation*}
on distribution~$\calD$. Then
\begin{equation*}
M_{a} \otimes I_{\reg{Bob}} \approx_{\delta} (\twirl_1 \otimes \mathrm{id}_{\mathrm{2}})[M_{a}] \otimes I_{\reg{Bob}}.
\end{equation*}
\end{lemma}
\begin{proof}
By definition,
\begin{equation*}
(\twirl_1 \otimes \mathrm{id}_{\mathrm{2}})[M_{a}]
= \E_{\bu, \bu'} [(X(\bu) Z(\bu') \otimes I_{\reg{2}}) \cdot M_a \cdot (Z(-\bu') X(-\bu) \otimes I_{\reg{2}})].
\end{equation*}
Similarly,
\begin{equation*}
M_{a}
= \E_{\bu, \bu'} [(X(\bu) Z(\bu') \otimes I_{\reg{2}}) \cdot (Z(-\bu') X(-\bu) \otimes I_{\reg{2}}) \cdot M_a].
\end{equation*}
As a result, if we set
$
A^{\bu, \bu'} = X(\bu) Z(\bu') \otimes I_{\reg{2}},
$
and
\begin{equation*}
B^{\bu, \bu'}_a = (Z(-\bu') X(-\bu) \otimes I_{\reg{2}}) \cdot M_a - M_a \cdot (Z(-\bu') X(-\bu) \otimes I_{\reg{2}}),
\end{equation*}
then
\begin{equation*}
\Delta_a:=
M_a - (\twirl_1 \otimes \mathrm{id}_{\mathrm{2}})[M_a]
= \E_{\bu, \bu'} [A^{\bu, \bu'} \cdot B^{\bu, \bu'}_a].
\end{equation*}
We can therefore establish the lemma as follows:
\begin{align*}
\sum_a \Vert (\Delta_a)_{\reg{Alice}} \otimes I_{\reg{Bob}} \ket{\psi} \Vert^2
& = \sum_a \Vert \E_{\bu, \bu'} [A^{\bu, \bu'} \cdot B^{\bu, \bu'}_a] \otimes I_{\reg{Bob}} \ket{\psi} \Vert^2\\
& \leq \E_{\bu, \bu'} \sum_a \Vert (A^{\bu, \bu'} \cdot B^{\bu, \bu'}_a) \otimes I_{\reg{Bob}} \ket{\psi} \Vert^2\tag{Jensen's inequality}\\
& = \E_{\bu, \bu'} \sum_a \Vert (B^{\bu, \bu'}_a) \otimes I_{\reg{Bob}} \ket{\psi} \Vert^2 \tag{$A^{\bu, \bu'}$ is unitary}
\end{align*}
By assumption, this quantity is $O(\delta)$. This concludes the proof.
\end{proof}


Now we prove \Cref{thm:data-hiding-game}.
\begin{proof}[Proof of \Cref{thm:data-hiding-game}]
We consider the completeness and soundness cases separately.

\paragraph{Completeness.}
Let $\calS_{\mathrm{partial}} = (\psi, M^x)$ be a partial
$(k,n,q)$-register strategy which is also a real commuting EPR strategy.
To this strategy we will add matrices for the questions $x' = (x_1', x_2)$ with $x_1' = (W_1, \ldots, W_{k-1}, W)$.

Let $a_1 = (u_1, \ldots, u_k)$, where $u_i = \varnothing$ if $W_i \in \{\hideq, \noop\}$.
Let~$S = \{i \mid W_i \neq \noop\}$.
By definition of a $(k, n, q)$-register strategy,
\begin{equation*}
M^x_a = \bigotimes_{i \in S} \tau^{W_i}_{u_i} \otimes M^{x, a_1}_{a_2},
\end{equation*}
where $M^{x,a_1}_{a_2}$ acts on the auxiliary registers and the registers not in~$S$.
Next, set $a' = (a_1', a_2)$ where $a_1' = (u_1, \ldots, u_{k-1}, u_k')$ and $u_k' \in \F_{q_k}^{n_k}$.  Then we set
\begin{equation*}
(M')^{x_1', x_2}_{a'_1, a_2} = \bigotimes_{i \in S\setminus k} \tau^{W_i}_{u_i} \otimes \tau^{W}_{u_k'} \otimes M^{x, a_1}_{a_2}.
\end{equation*}
This is a $(k,n,q)$-register strategy for~$\game_{\mathrm{hide}}$ by design.
To see that it is value~$1$, suppose on question~$x$ Player~$\bb$ measures~$\ba_1$.
Then by~\Cref{fact:the-ol-pauli-swaperoonie}, Player~$\overline{\bb}$ will measure~$\ba_1'$ in which $\bu_i' = \bu_i$ for all $i < k$.
As a result, to measure~$\ba_2$, Player~$\bb$ will measure $M^{x, \ba_1}$
and Player~$\overline{\bb}$ will measure $M^{x, \ba_1}$, both on state $\ket{r_{\overline{S}}}\ket{\mathrm{aux}}$.
As this is an EPR state,
by~\Cref{fact:heh-heh-heh-gonna-make-anand-prove-this-so-i-can-take-the-day-off}
the outcomes will always be the same, and so this strategy has
value~$1$. The fact that this a real strategy follows from the
assumption that the matrices $M^x_a$ are real, and the fact that for
$W \in \{X, Z\}$, $\tau^W_u$ is a real matrix.
Finally, the fact that this is a commuting strategy follows from the fact that $M^x_a$ and $(M')^{x'}_{a'}$ are commuting.

\paragraph{Soundness.}
We write $x$ and $x' = (x_1', x_2)$ with $x_1' = (W_1, \ldots, W_{k-1}, W)$ as in the test.
Because the test passes with probability $1-\eps$, \Cref{fact:agreement} implies that
\begin{equation*}
(M_{a_2}^{x})_{\reg{Alice}} \otimes I_{\reg{Bob}}
\approx_{\eps} I_{\reg{Alice}} \otimes (M_{a_2}^{x'})_{\reg{Bob}}.
\end{equation*}
Because $\calS$ is a $(k,n,q)$-semiregister strategy, 
\Cref{eq:in-math} implies that $M_{u_k}^{x'} = \tau^{W}_{u_k} \otimes I_{\overline{k}}$,
where we write $I_{\overline{k}} := I_{\reg{1, \ldots, k-1, aux}}$.
Our next step is to show that the measurements approximately commute.
This follows the analysis of the commutation test (cf.~\cite[Lemma 28]{CGJV18}).
\begin{align*}
M_{u_k}^{x'} M_{a_2}^{x} \otimes I_{\reg{Bob}}
&\approx_{\eps} M_{u_k}^{x'} \otimes M_{a_2}^{x'}\tag{\Cref{fact:add-a-proj}}\\
&\approx_{0} I_{\reg{Alice}} \otimes M_{a_2}^{x'}M_{u_k}^{x'} \tag{\Cref{fact:the-ol-pauli-swaperoonie}}\\
&= I_{\reg{Alice}} \otimes M_{u_k}^{x'}M_{a_2}^{x'}\\
&\approx_{\eps} M_{a_2}^{x} \otimes M_{u_k}^{x'}\tag{\Cref{fact:add-a-proj}}\\
&\approx_{0} M_{a_2}^x M_{u_k}^{x'} \otimes I_{\reg{Bob}}.\tag{\Cref{fact:the-ol-pauli-swaperoonie}}
\end{align*}
In summary,
\begin{equation*}
((\tau^{W}_{u_k} \otimes I_{\overline{k}})\cdot M_{a_2}^{x})_{\reg{Alice}} \otimes I_{\reg{Bob}}
\approx_{\eps} (M_{a_2}^x \cdot (\tau^{W}_{u_k} \otimes I_{\overline{k}}))_{\reg{Alice}}\otimes I_{\reg{Bob}}.
\end{equation*}
Recall this is with respect to the distribution $\bW$ where $\bW \sim \{X,Z\}$ is uniform.
Therefore, it also holds with respect to the distribution where $\bW$ is fixed to either~$X$ or~$Z$.
As a result, for a fixed $W \in \{X,Z\}$, by \Cref{lem:proj-to-obs}, 
\begin{equation*}
(M_{a_2}^x \cdot ( W(u)\otimes I_{\overline{k}}))_{\reg{Alice}} \otimes I_{\reg{Bob}}
\approx_{\eps} (( W(u)\otimes I_{\overline{k}}) \cdot M_{a_2}^x)_{\reg{Alice}} \otimes I_{\reg{Bob}}.
\end{equation*}
on distribution $\bu \sim \F_q^n$.
As a result, by~\Cref{fact:the-ol-pauli-swaperoonie} and \Cref{fact:add-a-proj},
\begin{align*}
(M_{a_2}^x \cdot (Z(u') X(u)\otimes I_{\overline{k}}))_{\reg{Alice}} \otimes I_{\reg{Bob}}
& \approx_{0}(M_{a_2}^x \cdot (Z(u')\otimes I_{\overline{k}}))_{\reg{Alice}} \otimes (X(-u)\otimes I_{\overline{k}})_{\reg{Bob}}\\
& \approx_{\eps}((Z(u')\otimes I_{\overline{k}}) \cdot M_{a_2}^x)_{\reg{Alice}} \otimes (X(-u)\otimes I_{\overline{k}})_{\reg{Bob}}\\
& \approx_{0}((Z(u')\otimes I_{\overline{k}}) \cdot M_{a_2}^x \cdot (X(u) \otimes I_{\overline{k}}))_{\reg{Alice}} \otimes I_{\reg{Bob}}\\
& \approx_{\eps}((Z(u') X(u)\otimes I_{\overline{k}}) \cdot M_{a_2}^x)_{\reg{Alice}} \otimes I_{\reg{Bob}},
\end{align*}
on distribution $\bu, \bu' \sim \F_q^n$.
Applying \Cref{lem:commute-to-twirl} and \Cref{prop:hidden-twirl},
we can therefore conclude
\begin{equation*}
M_{a_2}^x \otimes I_{\reg{Bob}}
\approx_{\eps} (\twirl_k \otimes \mathrm{id}_{\overline{k}})[M_{a_2}^x] \otimes I_{\reg{Bob}}
= (\hide{k}{M_{a_2}^x}) \otimes I_{\reg{Bob}}.\qedhere
\end{equation*}
\end{proof}



%%% Local Variables:
%%% mode: latex
%%% TeX-master: "main"
%%% End:
%%%---------- close: data-hiding

%%%%%%%%%%% jump to data-hiding-compiler
%%%---------- open: data-hiding-compiler
%!TEX root = main.tex

